% beamerposter 學術海報骨架 — A0 橫式三欄
% 編譯：pdflatex beamerposter_skeleton.tex（或 latexmk -pdf）
% 換主題改 \usetheme{...}；字太小調 beamerposter 的 scale
\documentclass[final]{beamer}
\usepackage[orientation=landscape,size=a0,scale=1.4]{beamerposter}
\usepackage{amsmath,amssymb}
\usepackage{graphicx}
\usetheme{Berlin}          % 可換：Madrid / Frankfurt / CambridgeUS ...

\title{My Research Poster Title}
\author{First Author \and Second Author}
\institute{Department, University}
\date{}

\begin{document}
\begin{frame}{}                         % 單一 frame = 整張海報
  \begin{columns}[t]                     % 頂端對齊的欄
    %------- 第 1 欄 -------
    \begin{column}{.32\linewidth}
      \begin{block}{Introduction}
        \begin{itemize}
          \item 研究動機 / 問題
          \item 關鍵相關工作
        \end{itemize}
      \end{block}
      \begin{block}{Methods}
        方法描述；公式原生渲染：
        \[ \mathcal{L} = \frac{1}{N}\sum_i \ell(y_i,\hat{y}_i) \]
      \end{block}
    \end{column}
    %------- 第 2 欄 -------
    \begin{column}{.32\linewidth}
      \begin{block}{Results}
        \begin{figure}
          \includegraphics[width=\linewidth]{results.pdf}  % 用向量 PDF
        \end{figure}
      \end{block}
    \end{column}
    %------- 第 3 欄 -------
    \begin{column}{.32\linewidth}
      \begin{block}{Conclusion}
        一兩行講清楚主要 takeaway。
      \end{block}
      \begin{block}{References}
        \footnotesize \bibliographystyle{plain}
      \end{block}
    \end{column}
  \end{columns}
\end{frame}
\end{document}
